\subsection{আইসিটি নির্ভর উৎপাদন ব্যবস্থা}

\begin{learningbox}
এই অংশ শেষে শিক্ষার্থী পারবে:
\begin{itemize}
\item ICT নির্ভর উৎপাদন ব্যবস্থার ধারণা ব্যাখ্যা করতে।
\item automation, computer-controlled machine, robot ও quality control-এর ভূমিকা লিখতে।
\item smart manufacturing ও Industry 4.0 ধারণার সঙ্গে ICT-এর সম্পর্ক বুঝতে।
\end{itemize}
\end{learningbox}

\subsubsection{স্বয়ংক্রিয় উৎপাদন}

উৎপাদন ব্যবস্থায় ICT ব্যবহার করলে machine, sensor, software, robot, database ও network একসঙ্গে কাজ করে। এতে raw material থেকে final product তৈরির ধাপ computer দ্বারা পরিকল্পিত, নিয়ন্ত্রিত ও পর্যবেক্ষিত হতে পারে।

\begin{definitionbox}{ICT নির্ভর উৎপাদন}
ICT নির্ভর উৎপাদন ব্যবস্থা হলো এমন manufacturing system, যেখানে computer, sensor, automation software, robot, database ও network ব্যবহার করে উৎপাদন, মান নিয়ন্ত্রণ, inventory ও delivery পরিচালনা করা হয়।
\end{definitionbox}

Smart manufacturing ধারণায় machine data, IoT, robotics, analytics ও automation ব্যবহার করে উৎপাদনকে দ্রুত, flexible ও efficient করা হয়।

\subsubsection{কম্পিউটার-নিয়ন্ত্রিত যন্ত্র}

CNC machine, 3D printer, automated packaging machine, PLC-controlled assembly line এবং industrial robot computer program অনুযায়ী কাজ করে। এতে shape, size, speed, temperature, pressure বা timing বেশি নির্ভুলভাবে control করা যায়।

\begin{examplebox}{CNC machine}
CNC machine computer program অনুসারে metal বা wood কাটে। একই design বারবার একই মাপে তৈরি করা যায়, তাই error কম হয়।
\end{examplebox}

\subsubsection{রোবট ব্যবহৃত উৎপাদন}

Robot welding, painting, material handling, assembling, packaging এবং inspection করতে পারে। মানুষের জন্য ঝুঁকিপূর্ণ, ভারী বা repetitive কাজ robot দিয়ে করালে নিরাপত্তা ও উৎপাদনশীলতা বাড়ে।

\subsubsection{গুণগত মান নিয়ন্ত্রণ}

Quality control-এ camera, sensor, barcode/RFID, database ও AI-based inspection ব্যবহার করা যায়। defective product শনাক্ত, production data record এবং batch tracking ICT দিয়ে করা যায়।

\begin{center}
\begin{tabular}{|p{0.28\textwidth}|p{0.58\textwidth}|}
\hline
\textbf{ICT tool} & \textbf{উৎপাদনে ভূমিকা} \\
\hline
Sensor & temperature, pressure, speed বা defect শনাক্ত \\
\hline
Database & inventory, batch, worker, order data সংরক্ষণ \\
\hline
Robot & repetitive ও hazardous কাজ করা \\
\hline
ERP software & finance, inventory, production ও sales সমন্বয় \\
\hline
Barcode/RFID & product tracking ও supply chain control \\
\hline
\end{tabular}
\end{center}

\subsubsection{উৎপাদন ব্যবস্থায় আইসিটির সুবিধা}

\begin{keypointbox}{সুবিধা}
\begin{itemize}
\item production speed বৃদ্ধি।
\item defect ও waste কমানো।
\item product quality consistent রাখা।
\item inventory ও supply chain tracking সহজ করা।
\item worker safety বৃদ্ধি।
\item উৎপাদন data বিশ্লেষণ করে planning উন্নত করা।
\end{itemize}
\end{keypointbox}

\begin{mistakebox}{সাধারণ ভুল}
Automation মানেই মানুষ অপ্রয়োজনীয় হয়ে যায়—এ ধারণা ঠিক নয়। Automation repetitive কাজ কমায়, কিন্তু design, maintenance, programming, supervision, quality analysis ও decision-making-এ দক্ষ মানুষের প্রয়োজন থাকে।
\end{mistakebox}

\begin{boardbox}{বোর্ড ফোকাস}
ICT নির্ভর উৎপাদনের উত্তরে automation, CNC, robot, sensor, database ও quality control শব্দগুলো ব্যবহার করলে উত্তর বেশি exam-ready হয়।
\end{boardbox}

\begin{practicebox}
\textbf{MCQ}
\begin{enumerate}
\item CNC machine কোনটির দ্বারা নিয়ন্ত্রিত হয়?\\
ক. computer program \quad খ. শুধুই হাতুড়ি \quad গ. কাগজের খাতা \quad ঘ. radio only
\item Barcode/RFID উৎপাদনে কোন কাজে সহায়তা করে?\\
ক. product tracking \quad খ. poem writing \quad গ. screen brightness \quad ঘ. keyboard cleaning
\end{enumerate}

\textbf{সংক্ষিপ্ত প্রশ্ন}
\begin{enumerate}
\item ICT নির্ভর উৎপাদন ব্যবস্থা কী?
\item উৎপাদনে sensor ব্যবহারের দুটি সুবিধা লেখ।
\end{enumerate}
\end{practicebox}

\begin{sourcebox}
এই অংশে local ICT PDF resources, NIST smart manufacturing ধারণা, Industry 4.0 articles এবং manufacturing automation examples মিলিয়ে লেখা হয়েছে।
\end{sourcebox}
